\documentclass[11pt]{article}
\usepackage{amsmath,amssymb,amsthm}
\usepackage{enumitem}
\usepackage[margin=1.25in]{geometry}
\usepackage[colorlinks=true,linkcolor=black,citecolor=black,urlcolor=black]{hyperref}

% -------------------- macros --------------------
\newcommand{\E}{\mathbb{E}}
\newcommand{\Prob}{\mathbb{P}}
\newcommand{\R}{\mathbb{R}}
\newcommand{\N}{\mathbb{N}}
\newcommand{\1}{\mathbf{1}}
\newcommand{\Var}{\mathrm{Var}}
\newcommand{\Cov}{\mathrm{Cov}}
\newcommand{\Tr}{\mathrm{Tr}}
\newcommand{\eps}{\varepsilon}
\newcommand{\cyc}{\mathrm{cyc}}
\newcommand{\pd}{\mathrm{PD}(1)}
\newcommand{\PD}{\mathrm{PD}}
\newcommand{\GEM}{\mathrm{GEM}(1)}
\newcommand{\Poisson}{\mathrm{Poisson}}
\newcommand{\gE}{\gamma_{\mathrm{E}}}
\newcommand{\Err}{\mathrm{Err}}
\newcommand{\LIS}{\mathrm{LIS}}
\newcommand{\cA}{\mathcal{A}}
\newcommand{\cB}{\mathcal{B}}
\newcommand{\cE}{\mathcal{E}}
\newcommand{\cF}{\mathcal{F}}
\newcommand{\cG}{\mathcal{G}}
\newcommand{\cH}{\mathcal{H}}
\newcommand{\cI}{\mathcal{I}}
\newcommand{\cN}{\mathcal{N}}
\newcommand{\cS}{\mathcal{S}}
\newcommand{\cT}{\mathcal{T}}

\newtheorem{theorem}{Theorem}[section]
\newtheorem{lemma}[theorem]{Lemma}
\newtheorem{proposition}[theorem]{Proposition}
\newtheorem{corollary}[theorem]{Corollary}
\newtheorem{estimate}[theorem]{Estimate}
\theoremstyle{remark}
\newtheorem{remark}[theorem]{Remark}

\title{Mesoscopic cycles of Luce permutations:\\
resolution of Conjectures 13.1 and 13.2 of \cite{Lon26}}
\author{}
\date{}

\begin{document}
\maketitle

\begin{abstract}
We resolve the mesoscopic cycle conjectures of the revised manuscript
\cite{Lon26} on Luce permutations with comparable weights. Under the
comparable-weights hypothesis \eqref{eq:comparable} we prove the local
cycle-length law $\Prob(L_n=r)=(1+o(1))/n$, uniformly for
$r_n\le r\le\delta_0 n$ whenever $r_n\to\infty$ (Theorem~\ref{thm:local}),
with an explicit rate. Its consequences include
$\E C_{r_n}(\sigma_n)\sim 1/r_n$, a joint local law for any fixed
number of lengths, and convergence of the rescaled mesoscopic cycle
lengths to the scale-invariant Poisson process with intensity $dx/x$,
which is Conjecture~13.1 of \cite{Lon26} in full. Under the
sampled-profile hypothesis \eqref{eq:profile} we then prove
$\E N_n^{\cyc}=\log n+\kappa_f+o(1)$ with the conjectured constant
$\kappa_f$, and the central limit theorem
$(N_n^{\cyc}-\E N_n^{\cyc})/\sqrt{\log n}\Rightarrow N(0,1)$, which is
Conjecture~13.2 of \cite{Lon26} in full. For the remaining
conjectures we give precise reductions and identify the exact
obstructions: Conjecture~A.1 follows from a stated mixed-scale local
estimate (Theorem~\ref{thm:A1-reduction}), Conjecture~13.3 reduces to
the local law at the endpoint scale, where the comparable-weights
hypothesis degenerates (Proposition~\ref{prop:133}), and
Conjecture~A.2 requires control at scale $\Theta(n/\log n)$ that is
beyond the present method. All probabilistic inputs from \cite{Lon26}
are quoted with their precise statement numbers; the new dynamical
estimates are proved in full. The expectation constant and the
mesoscopic scale-invariance are verified numerically in
Appendix~\ref{app:numerics}.
\end{abstract}

\section{Introduction and statement of results}\label{sec:intro}

\subsection{The model}
Let $X_{n,i}\sim\mathrm{Exp}(\theta_{n,i})$, $1\le i\le n$, be
independent clocks with rates $\theta_{n,i}>0$, and let $\sigma_n$ be
the label-to-position \emph{Luce permutation}
\begin{equation*}
\sigma_n(i)\;=\;1+\#\{j:X_{n,j}<X_{n,i}\},
\end{equation*}
so that $\sigma_n(i)=k$ when $X_{n,i}$ is the $k$-th smallest clock
(the ranking model of \cite{Luc59}; see \cite[Section~2]{Lon26} and
\cite{CDK24,BCD25} for background). Two hypotheses are used in this
paper, always exactly as in \cite{Lon26}. The \emph{comparable-weights}
hypothesis, underlying \cite[Theorem~6.1]{Lon26}: since multiplying
all rates by a constant does not change $\sigma_n$, we may normalize
so that for some fixed $\kappa<\infty$,
\begin{equation}\label{eq:comparable}
1\;\le\;\theta_{n,i}\;\le\;\kappa\qquad\text{for all }n,i.
\end{equation}
The \emph{exact sampled-profile} hypothesis, underlying
\cite[Section~5]{Lon26} and Conjecture~13.2 there: for some
$f\in C[0,1]$ with $0<m_f\le f\le M_f<\infty$,
\begin{equation}\label{eq:profile}
\theta_{n,i}\;=\;f\bigl(i/n\bigr),\qquad 1\le i\le n.
\end{equation}
Under \eqref{eq:profile}, with $A(t)=\int_0^1(1-e^{-f(u)t})\,du$,
$B=A'$ and $\tau=A^{-1}$, the limiting local transition kernel (the
Luce permuton density of \cite[Proposition~2.2]{BCD25}, quoted as
\cite[Section~5]{Lon26}) is
\begin{equation}\label{eq:permuton}
\rho(x,y)\;=\;\frac{f(x)\,e^{-f(x)\tau(y)}}{B(\tau(y))},
\qquad 0\le x\le1,\ 0<y<1,
\end{equation}
and $K$ denotes the bistochastic integral operator
$(Kh)(x)=\int_0^1\rho(x,y)h(y)\,dy$, with eigenvalues
$1=\mu_1>|\mu_2|\ge\cdots$ and traces
$r\lambda_r=\Tr(K^r)$ \cite[Theorem~5.19]{Lon26}.

Write $C_r(\sigma_n)$ for the number of cycles of $\sigma_n$ of length
$r$, $N_n^{\cyc}=\sum_{r\ge1}C_r(\sigma_n)$ for the total number of
cycles, $L_n(v)$ for the length of the cycle containing $v\in[n]$,
$L_{n,1}\ge L_{n,2}\ge\cdots$ for the ranked cycle lengths, and $I_n$
for a uniform element of $[n]$ independent of the clocks.

\subsection{Results: Conjectures 13.1 and 13.2}
Conjecture~13.1 of \cite{Lon26} predicts, under the hypotheses of
\cite[Theorem~6.1]{Lon26} (i.e.\ \eqref{eq:comparable}), that
$\E C_{r_n}(\sigma_n)\sim 1/r_n$ whenever $r_n\to\infty$ with
$r_n=o(n)$, and more strongly that for $a_n\to\infty$, $a_n=o(n)$ and
disjoint $[u_j,v_j]\Subset(0,\infty)$, the counts
$\#\{i:u_ja_n\le L_{n,i}\le v_ja_n\}$ converge to independent Poisson
variables with means $\log(v_j/u_j)$; equivalently the rescaled
submacroscopic cycle lengths converge to the Poisson process with
intensity $dx/x$. Conjecture~13.2 of \cite{Lon26} predicts, under
\eqref{eq:profile}, that $\E N_n^{\cyc}=\log n+\kappa_f+o(1)$ with
$\kappa_f=\gE+\sum_{r\ge1}(\lambda_r-\frac1r)$ (the series converges
absolutely by \cite[Theorem~5.19]{Lon26}), and that
$(N_n^{\cyc}-\E N_n^{\cyc})/\sqrt{\log n}\Rightarrow N(0,1)$.

Both conjectures follow from a single new estimate, a local law for
the cycle length of a uniform root, uniform over the entire mesoscopic
range.

\begin{theorem}[Local cycle-length law]\label{thm:local}
Assume \eqref{eq:comparable}. There exist $\delta_0=\delta_0(\kappa)>0$
and constants $c_1,c_2\in(0,1)$ depending only on $\kappa$ such that,
for every sequence $r_n\to\infty$,
\begin{equation}\label{eq:local-law}
\Prob\bigl(L_n(I_n)=r\bigr)
\;=\;\frac{1}{n}\Bigl(1+O\bigl(c_1^{\,r}\bigr)+O\bigl(n^{-c_2}\bigr)\Bigr)
\;=\;\frac{1+o(1)}{n},
\end{equation}
uniformly in $r_n\le r\le\delta_0 n$; the $o(1)$ is uniform in the
range.
\end{theorem}

\begin{corollary}[Expected mesoscopic cycle counts]\label{cor:EC}
Assume \eqref{eq:comparable}. If $r_n\to\infty$ with $r_n=o(n)$, then
\[
\E C_{r_n}(\sigma_n)\;=\;\frac{1+o(1)}{r_n}.
\]
This is the first part of \cite[Conjecture~13.1]{Lon26}.
\end{corollary}

\begin{theorem}[Joint local law]\label{thm:joint-local}
Assume \eqref{eq:comparable}. Fix $\ell\ge1$ and let
$V_1,\dots,V_\ell$ be independent uniform elements of $[n]$,
independent of $\sigma_n$. For any lengths
$k_{n,1},\dots,k_{n,\ell}\to\infty$ with
$\sum_{j=1}^\ell k_{n,j}=o(n)$,
\begin{equation}\label{eq:joint-local}
\Prob\Bigl(L_n(V_j)=k_{n,j}\ \text{for all }j\le\ell,\ \text{and the
cycles of }V_1,\dots,V_\ell\text{ are distinct}\Bigr)
\;=\;\frac{1+o(1)}{n^\ell},
\end{equation}
with $o(1)$ uniform in the lengths.
\end{theorem}

\begin{theorem}[Mesoscopic Poisson process]\label{thm:ppp}
Assume \eqref{eq:comparable}. Let $a_n\to\infty$ with $a_n=o(n)$ and
let $[u_j,v_j]\Subset(0,\infty)$, $1\le j\le m$, be pairwise disjoint.
Then
\begin{equation}\label{eq:ppp-vector}
\Bigl(\#\bigl\{i:u_ja_n\le L_{n,i}\le v_ja_n\bigr\}\Bigr)_{j=1}^m
\;\Rightarrow\;\bigl(Z_j\bigr)_{j=1}^m,
\end{equation}
where $Z_j$ are independent with
$Z_j\sim\Poisson\bigl(\log(v_j/u_j)\bigr)$. Equivalently, the point
process $\sum_{i}\delta_{L_{n,i}/a_n}$ converges, on compact subsets
of $(0,\infty)$, to the Poisson process with intensity $dx/x$. This
resolves \cite[Conjecture~13.1]{Lon26} in full.
\end{theorem}

\begin{theorem}[Expected number of cycles]\label{thm:expectation}
Assume \eqref{eq:profile}. Then
\[
\E N_n^{\cyc}\;=\;\log n+\kappa_f+o(1),
\qquad
\kappa_f\;=\;\gE+\sum_{r=1}^\infty\Bigl(\lambda_r-\frac1r\Bigr),
\]
with $\lambda_r=\Tr(K^r)/r$; the series converges absolutely by
\cite[Theorem~5.19]{Lon26}. Equivalently
$\kappa_f=\gE-\sum_{i\ge2}\log(1-\mu_i)$, where $1=\mu_1,\mu_2,\dots$
are the eigenvalues of $K$. This is the first part of
\cite[Conjecture~13.2]{Lon26}.
\end{theorem}

\begin{theorem}[Central limit theorem]\label{thm:clt}
Assume \eqref{eq:profile}. Then
\[
\frac{N_n^{\cyc}-\E N_n^{\cyc}}{\sqrt{\log n}}\;\Rightarrow\;N(0,1),
\qquad\text{and}\qquad
\Var N_n^{\cyc}\;=\;\log n+o(\log n).
\]
Together with Theorem~\ref{thm:expectation} this resolves
\cite[Conjecture~13.2]{Lon26} in full.
\end{theorem}

\begin{remark}[Uniformity and range]\label{rem:uniform}
The local law \eqref{eq:local-law} is uniform for
$r_n\le r\le\delta_0 n$; Conjecture~13.1 only asks for $r_n=o(n)$, so
the upper end of our range ($r$ a small positive proportion of $n$)
exceeds what is required, and the lower end ($r_n\to\infty$
arbitrarily slowly, e.g.\ $r_n=\log\log n$) is where the difficulty
lies: the proof must control orbit mixing at scales at which the
spectral gap has had only $O(\log\log n)$ steps to act. The rate
$O(c_1^{\,r})+O(n^{-c_2})$ in \eqref{eq:local-law} is used essentially
in the expectation computation, where the local law is summed over
$\Theta(n)$ scales.
\end{remark}

\subsection{Overview of the proof}
The proof of Theorem~\ref{thm:local} is built on the clock-slab
framework of \cite[Section~6]{Lon26}, recalled in
Section~\ref{sec:prelim} with exact citations: the row decomposition
of the permutation into $\cH_n$-measurable rank intervals, the
within-row permanental matchings, the hereditary flatness of one-edge
marginals, the exactly stationary compressed kernel $P_n$ with its
spectral gap on the good event $\cG_n$, and the root-avoiding change
of measure $\Prob^\circ$ of \cite[Lemma~6.9]{Lon26}, under which
closure to the root can occur only at visits of the orbit to the
root's rank row. The new ingredients are:
\begin{enumerate}[label=(\alph*)]
\item \emph{Balanced depletion} (Lemma~\ref{lem:balance}): a
quantitative form of the relative-prefix device of
\cite[Lemma~6.7]{Lon26}, proved by Freedman's martingale inequality,
showing that after $r\le\delta n$ exposures every cell of every row is
depleted proportionally to its size, up to a polynomially small
relative error;
\item \emph{Mixing of the orbit's relative density}
(Lemmas~\ref{lem:onestep} and \ref{lem:mixing}): the one-step slab
kernel is $(1+o(1))P_n$ uniformly over histories, hence the density
$u_k$ of the orbit's slab law relative to the stationary law $p_n$
contracts to $1$ at the geometric rate of the gap, initialized after a
single step via the ratio bound \eqref{eq:ratio-bound};
\item \emph{Visit-count fluctuations} (Lemma~\ref{lem:visits}): the
number of closure opportunities $S(k)$ concentrates at
$kp_{n,c(v)}=O(\delta nh)$, with a covariance bound inherited from the
gap;
\item \emph{The root average in trace form}
(Lemma~\ref{lem:root-avg}): averaging the return density over a
uniform root gives exactly an expected trace,
$R_k=\E^\circ\Tr(P_n\tilde P^{(1)}\cdots\tilde P^{(k)})$, the
finite-volume shadow of the trace universality
$r\lambda_r=\Tr(K^r)=1+O(q^r)$ of \cite[Theorem~5.19]{Lon26}; a
perturbation expansion using the exact identities $p_nQ=0$, $Q1=0$
controls it for $k\le c_*\log n$ (the scale-adaptive mesh), and a
half-split Hilbert--Schmidt bound controls it beyond (the fixed fine
mesh);
\item \emph{Hazard telescoping and the two mesh regimes}
(Lemmas~\ref{lem:hazard} and \ref{lem:telescope},
Remark~\ref{rem:two-mesh}): the closure probability factors into an
exact telescoping product, $(1+O(\beta))/n$ times the return density,
with $\beta=2\kappa h$; the scale-adaptive mesh $h_k=(q')^k$ makes all
errors vanish for slowly growing $k$, while the fixed mesh
$h=n^{-1/8}$ gives a polynomially small error for
$k\asymp\log n$ and beyond.
\end{enumerate}
Theorem~\ref{thm:joint-local} then follows by iterated residual
exploration (Section~\ref{sec:joint}): deleting completely explored
cycles leaves square rows, a $(1+o(1))$-perturbed kernel with the same
gap, and the exact cancellation $(1-K/n)/(n-K)^{-1}=1/n$.
Theorem~\ref{thm:ppp} follows from the joint factorial moments of the
window counts by the method of moments (Section~\ref{sec:ppp}).
Theorem~\ref{thm:expectation} is proved in
Section~\ref{sec:expectation} by a four-range decomposition in which
the fixed-length traces $\lambda_r$ of \cite[Theorem~5.1]{Lon26}, the
harmonic mesoscopic sum, and the macroscopic count
$\log(1/\delta)$ of \cite[Proposition~6.14]{Lon26} meet with exact
cancellation of $\log\delta$; the constant $\kappa_f$ is the
accumulation of the trace corrections $\lambda_r-1/r$ across all
scales. Theorem~\ref{thm:clt} is proved in Section~\ref{sec:clt} by
computing the factorial moments of the mesoscopic cycle count and
applying the method of moments, the short and macroscopic parts being
tight at scale $\sqrt{\log n}$.

\subsection{The remaining conjectures}
Section~\ref{sec:remaining} treats the conjectures we do not resolve.
Conjecture~A.1 of \cite{Lon26} (bulk--endpoint independence) is
reduced, completely and explicitly, to a mixed-scale local estimate
(Estimate~\ref{est:mixed}, Theorem~\ref{thm:A1-reduction}) whose two
extreme marginals are Theorem~\ref{thm:joint-local} here and
\cite[Theorem~10.1]{Lon26}; the missing ingredient is an exclusion
argument for orbits of length comparable to $n$
(Remark~\ref{rem:est-status}). Conjecture~13.3 (endpoint-singular
$\PD(1)$) is reduced to the local law at the endpoint scale
$k=\Theta(n)$ (Proposition~\ref{prop:133}); the obstruction is that
the weight ratio $\kappa_n\to\infty$ degenerates the spectral gap, the
flatness window, and the ratio bound simultaneously
(Remark~\ref{rem:obstruction}). Conjecture~A.2 requires the cycle
upper tail at scale $\Theta(n/\log n)$, where the present local law
ceases to be effective (Section~\ref{sec:A2}).

\subsection{Numerical verification}
Appendix~\ref{app:numerics} reports direct simulations for the exact
profile $f(x)=1+x$: the spectral data of $K$ give
$\kappa_f=0.4410$, and the simulated $\E N_n^{\cyc}$ at
$n=500,1000,2000,4000$ match $\log n+\kappa_f$ to two decimals; the
dyadic mesoscopic windows have mean $\approx\log2$ as predicted by
Theorem~\ref{thm:ppp}; and $\Var N_n^{\cyc}/\log n$ increases toward
$1$ as predicted by Theorem~\ref{thm:clt}.

\subsection{Conventions}
References of the form ``\cite[Lemma~6.9]{Lon26}'' cite statement
numbers of the revised manuscript \cite{Lon26}; all other numbered
statements are ours. Every input from \cite{Lon26} is quoted by number
where it is used, and everything else is proved. The letters $C,c$
denote finite positive constants depending only on $\kappa$ (and on
$f$ under \eqref{eq:profile}) whose value may change from line to
line; $q_0,q',c_1,c_2\in(0,1)$ are likewise $\kappa$-dependent.


\section{Preliminaries: the clock-slab framework of \cite{Lon26}}\label{sec:prelim}

This section recalls the framework of \cite[Section~6]{Lon26} in the form we
use it, with every input quoted by its precise statement number. Nothing in
this section is new.

\subsection{Slabs, the row decomposition, and the compressed kernel}
Let $0=t_0<t_1<\cdots<t_d=T$ be a partition of $[0,T]$ (the \emph{master
partition}) with mesh $\max_a(t_a-t_{a-1})\le h$, where $T=T(n)$ and the mesh
$h=h(n)$ are chosen later; $h$ is the only free parameter of the framework
and will be tuned per regime (Section~\ref{sec:orbit}). Write
$I_a=(t_{a-1},t_a]$ and
\[
H_{n,a}\;=\;\{i\in[n]:X_{n,i}\in I_a\},\qquad m_{n,a}\;=\;|H_{n,a}|,\qquad
p_{n,a}\;=\;\frac{m_{n,a}}{n},\qquad a\in[d].
\]
Let $\cH_n$ be the $\sigma$-field generated by the complete slab data (the
sets $H_{n,a}$ and the ordered clocks inside them). The following structural
facts are established in \cite[Section~6.2]{Lon26}.

\begin{enumerate}[label=(S\arabic*)]
\item \textbf{Rank intervals.} \cite[Section~6.2]{Lon26}: there are
$\cH_n$-measurable rank intervals
$Q_{n,a}=\{M_{n,a-1}+1,\dots,M_{n,a}\}$ partitioning $[n]$ with
$|Q_{n,a}|=m_{n,a}$, such that the permutation respects the rows: writing
$c(v)$ for the unique $a$ with $v\in Q_{n,a}$ (the \emph{rank row} of $v$)
and $b(v)$ for the unique $b$ with $v\in H_{n,b}$ (the \emph{clock slab} of
$v$), one has: $\sigma_n$ maps $C\cap H_{n,a}$ bijectively onto
$C\cap Q_{n,a}$ for every cycle $C$ of $\sigma_n$; in particular each orbit
step maps a source in $H_{n,a}$ to a target in $Q_{n,a}$.
\item \textbf{Within-row permanental matchings.}
\cite[Lemma~6.4]{Lon26}: conditional on $\cH_n$ and on any compatible partial
collection of exposed matching edges, the unexposed part of the matching in
each row $a$ (sources $H_{n,a}$, targets $Q_{n,a}$) is a permanental matching
on its remaining submatrix, independently across rows.
\item \textbf{Flatness.} \cite[Lemmas~6.5, 6.6]{Lon26}: every one-point
weight ratio within a row of width $\le h$ lies in $[e^{-\beta},e^{\beta}]$
with $\beta=2\kappa h$, and this flatness is hereditary under exposures:
conditional on any compatible history, if a row has $s$ residual targets then
every residual target has conditional probability in $[e^{-\beta}/s,
e^{\beta}/s]$. The same holds with $s-1$ in place of $s$ if one distinguished
target is forbidden \cite[Lemma~6.6]{Lon26}.
\item \textbf{Compressed kernel.} \cite[Section~6.2]{Lon26}: with
\[
N_{n,ab}\;=\;|Q_{n,a}\cap H_{n,b}|,\qquad
P_{n,ab}\;=\;\frac{N_{n,ab}}{m_{n,a}},
\]
the matrix $P_n=(P_{n,ab})_{a,b\in[d]}$ is a Markov kernel on $[d]$
($\sum_bN_{n,ab}=|Q_{n,a}|=m_{n,a}$) with stationary distribution $p_n$:
indeed $\sum_a p_{n,a}P_{n,ab}=\frac1n\sum_a N_{n,ab}=\frac{|H_{n,b}|}{n}
=p_{n,b}$. Under \eqref{eq:comparable} there is $A=A(\kappa)<\infty$ with
\begin{equation}\label{eq:ratio-bound}
\kappa^{-A}\;\le\;\frac{P_{n,ab}}{p_{n,b}}\;\le\;\kappa^{A}
\qquad\text{for all }a,b,
\end{equation}
by the within-row flatness and the comparability of the clock densities
(each $X_{n,i}$ has density in $[\kappa^{-1}e^{-\kappa T},1]$ on $[0,T]$).
\end{enumerate}

The orbit of a root $v$ is $q_0=v$, $q_{k+1}=\sigma_n(q_k)$; its slab chain
is $b_k=b(q_k)$. Since $\sigma_n(q_k)\in Q_{n,b_k}$ and the next source
$q_{k+1}=\sigma_n(q_k)$ is a flat draw among the residual targets
of row $b_k$ (S3),
\begin{equation}\label{eq:slab-transition}
\Prob\bigl(b_{k+1}=b\mid b_k=a,\ \text{history}\bigr)
\;=\;e^{O(\beta)}\,\frac{\widetilde N_{ab}}{\widetilde m_a},
\end{equation}
where $\widetilde N_{ab}=|Q_{n,a}\cap H_{n,b}\cap\{\text{residual}\}|$ and
$\widetilde m_a=\sum_b\widetilde N_{ab}$ are the residual cell and row
counts. At the start of the exploration $\widetilde N_{ab}=N_{n,ab}$ and
\eqref{eq:slab-transition} equals $e^{O(\beta)}P_{n,ab}$; controlling the
exploration means controlling the ratio $\widetilde N_{ab}/\widetilde m_a$
(Lemmas~\ref{lem:balance} and \ref{lem:onestep}).

\begin{proposition}[Good event and spectral gap; \cite{Lon26}]\label{prop:good}
Under \eqref{eq:comparable}, provided $nh\to\infty$ polynomially:
\begin{enumerate}[label=(\roman*)]
\item \textbf{Good event} \cite[Lemma~6.3]{Lon26}: there is an
$\cH_n$-measurable event $\cG_n$ on which: all $m_{n,a}$ and $N_{n,ab}$
concentrate at their means at the natural fluctuation scale; in particular
\begin{equation}\label{eq:pmin}
m_{n,a}\;\ge\;cnh\quad\text{for all }a,
\qquad\text{i.e.}\qquad p_{n,a}\;\ge\;ch,
\end{equation}
since each clock has density $\ge\kappa^{-1}e^{-\kappa T}$ on $[0,T]$; and
$H_{n,\mathrm{tail}}=\{i:X_{n,i}>T\}$ is empty. For $T=(B+1)\log n$ one has
$\Prob(\cG_n^c)=O(n^{-B})$ for the prescribed $B<\infty$, the constant
depending on $B,\kappa$ and the deviation parameter of
\cite[Lemma~6.3]{Lon26}.
\item \textbf{Spectral gap} \cite[Lemmas~6.2, 6.3]{Lon26}: there is
$q_0=q_0(\kappa)\in(0,1)$ such that on $\cG_n$,
\begin{equation}\label{eq:gap}
\|P_n v\|_{L^2(p_n)}\;\le\;q_0\,\|v\|_{L^2(p_n)}
\qquad\text{whenever }\E_{p_n}[v]=0,
\end{equation}
and the same bound holds for the adjoint $P_n^*$ with respect to $p_n$.
\end{enumerate}
\end{proposition}

The number of slabs is $d=\Theta(T/h)=O(h^{-1}\log n)$, so
\begin{equation}\label{eq:d-bound}
\sum_a p_{n,a}^{1/2}\;\le\;d^{1/2}\;=\;O\bigl(h^{-1/2}\sqrt{\log n}\bigr)
\end{equation}
by Cauchy--Schwarz.

\begin{proposition}[Root-avoiding change of measure; \cite{Lon26}]\label{prop:rootavoid}
Fix a root $v$ and condition on $\cH_n$. Explore the orbit from $v$, exposing
the edge $\sigma_n(q_t)=q_{t+1}$ at step $t$. Let $\Prob^\circ$ be the
auxiliary law in which, at every visit of the orbit to the root's rank row
$c(v)$, the within-row selection is conditioned to avoid the target $v$
\cite[Lemma~6.9]{Lon26}. Then
\begin{equation}\label{eq:survival}
\Prob\bigl(L_n(v)>k\mid\cH_n\bigr)
\;=\;\E^\circ\Bigl[\prod_{j=0}^{S(k)-1}(1-\lambda_j)\Bigr],
\end{equation}
where $S(k)=\#\{0\le t\le k-1:b_t=c(v)\}$ is the number of closure
opportunities among the first $k$ exposed edges and, on $\cG_n$, the closure
marginals satisfy
\begin{equation}\label{eq:hazard-bounds}
\frac{e^{-\beta}}{m_{n,c(v)}-j}\;\le\;\lambda_j\;\le\;
\frac{e^{\beta}}{m_{n,c(v)}-j},\qquad \beta=2\kappa h,
\end{equation}
by \cite[Lemmas~6.6, 6.9]{Lon26}. Under $\Prob^\circ$ the orbit evolves by the
residual permanental dynamics with target $v$ forbidden, to which the flat
one-edge marginals (S3) still apply \cite[Lemma~6.6]{Lon26}.
\end{proposition}

\begin{proposition}[Macroscopic limit; \cite{Lon26}]\label{prop:macro}
Under \eqref{eq:comparable}, the ranked macroscopic cycle lengths converge to
$\pd$ in $S^\downarrow$ and the size-biased ones to $\GEM$ in $\ell^1$
\cite[Theorem~6.1]{Lon26}; for every fixed $\delta>0$,
$\#\{i:L_{n,i}\ge\delta n\}$ converges in distribution to
$\Poisson(\log(1/\delta))$ with convergence of all moments
\cite[Proposition~6.14]{Lon26}; in particular
\begin{equation}\label{eq:macro-count}
\E\,\#\{i:L_{n,i}\ge\delta n\}\;\longrightarrow\;\log\frac1\delta
\qquad(n\to\infty,\ \delta>0\text{ fixed}).
\end{equation}
\end{proposition}

\subsection{Short cycles and traces under the sampled profile}
Under \eqref{eq:profile} we cite:

\begin{proposition}[Short cycles and traces; \cite{Lon26}]\label{prop:short}
Under \eqref{eq:profile}:
\begin{enumerate}[label=(\roman*)]
\item \textbf{Fixed cycle lengths} \cite[Theorem~5.1]{Lon26}: for every fixed
$R$, $(C_1(\sigma_n),\dots,C_R(\sigma_n))\Rightarrow(Z_1,\dots,Z_R)$ with
$Z_r$ independent $\Poisson(\lambda_r)$; the proof in \cite{Lon26} is by the
method of factorial moments and yields
\begin{equation}\label{eq:fixed-fm}
\E\prod_{r=1}^R\bigl(C_r(\sigma_n)\bigr)_{k_r}\;\longrightarrow\;
\prod_{r=1}^R\lambda_r^{k_r}\qquad\text{for all fixed }k_1,\dots,k_R,
\end{equation}
where $(x)_k=x(x-1)\cdots(x-k+1)$; together with the domination bounds of
\cite[Lemmas~5.8--5.13]{Lon26} this gives $\E C_r(\sigma_n)\to\lambda_r$ for
each fixed $r$.
\item \textbf{Separated local estimate} \cite[Proposition~5.7]{Lon26}: for
separated bulk constraints $(p_a,q_a)$, $a\le\ell$ fixed,
\begin{equation}\label{eq:separated}
\Prob\bigl(\sigma_n(p_a)=q_a\ \forall a\le\ell\bigr)
\;=\;\frac{1+o(1)}{n^\ell}\prod_{a=1}^\ell \rho\bigl(p_a/n,q_a/n\bigr).
\end{equation}
\item \textbf{Trace universality} \cite[Theorem~5.19, Lemmas~5.17--5.18]
{Lon26}: the operator $K$ with kernel \eqref{eq:permuton} is bistochastic and
Hilbert--Schmidt, with eigenvalues $1=\mu_1>\lvert\mu_2\rvert\ge\cdots$
satisfying $\lvert\mu_i\rvert\le q$ for $i\ge2$ and some $q=q(f)\in(0,1)$;
with $\lambda_r=\Tr(K^r)/r$,
\begin{equation}\label{eq:trace-gap}
r\lambda_r\;=\;1+O(q^r),\qquad
\sum_{r=1}^\infty\Bigl|\lambda_r-\frac1r\Bigr|<\infty,\qquad
\kappa_f\;=\;\gE-\sum_{i\ge2}\log(1-\mu_i).
\end{equation}
\end{enumerate}
\end{proposition}

The rank-envelope bounds \cite[Lemmas~5.8--5.10]{Lon26} and the endpoint
domination bounds \cite[Lemmas~5.11--5.13]{Lon26} are used only through item
(i) above and in Section~\ref{sec:remaining}.

\subsection{Standing notation}
Throughout, $C,c$ denote finite positive constants depending only on $\kappa$
(and on $f$ under \eqref{eq:profile}), whose value may change from line to
line; $q_0,q',c_1,c_2\in(0,1)$ likewise. We write
$p_{n,\min}=\min_a p_{n,a}$, so $p_{n,\min}\ge ch$ on $\cG_n$ by
\eqref{eq:pmin}. For a root $v$ we write $b(v)$ for its clock slab and $c(v)$
for its rank row as in (S1); the orbit is $q_0=v$, $q_{k+1}=\sigma_n(q_k)$
with slab chain $b_k=b(q_k)$; $\Prob^\circ$ is the root-avoiding law of
Proposition~\ref{prop:rootavoid}. The \emph{relative density} of the orbit's
slab law with respect to $p_n$ is
\[
u_k(b)\;=\;\frac{\Prob^\circ(b_k=b\mid\cH_n)}{p_{n,b}},\qquad b\in[d],
\]
which is $\cH_n$-conditionally a probability density: $\E_{p_n}[u_k]=1$.
All statements below are conditional on $\cH_n$ on the event $\cG_n$ unless
explicitly averaged. Finally, $q_0$ from \eqref{eq:gap} is a
$\kappa$-dependent bound, and a standard trace bound follows from
\eqref{eq:ratio-bound}: since $P_{n,ab}\le\kappa^A p_{n,b}$,
\begin{equation}\label{eq:hs}
\|P_n\|_{\mathrm{HS}}^2\;=\;\sum_{a,b}P_{n,ab}^2
\;\le\;\kappa^A\sum_{a,b}p_{n,b}P_{n,ab}\;=\;\kappa^A,
\end{equation}
and any product $M=P_nM_1\cdots M_s$ of length $\ge1$ with Markov factors
satisfies the same ratio bound $M_{ab}\le\kappa^Ap_{n,b}$, hence
$\|M\|_{\mathrm{HS}}^2\le\kappa^A$ and $\|M\|_{\mathrm{tr}}\le
d^{1/2}\kappa^{A/2}$.

\section{Orbit dynamics in slab space}\label{sec:orbit}

This section contains the four new dynamical estimates: a balanced-depletion
lemma (Lemma~\ref{lem:balance}), the one-step transition estimate
(Lemma~\ref{lem:onestep}), the uniform mixing estimate
(Lemma~\ref{lem:mixing}), and the trace-form root average
(Lemma~\ref{lem:root-avg}).

\subsection{Balanced depletion}
The next lemma is a quantitative form of the relative-prefix device of
\cite[Lemma~6.7]{Lon26}, with failure probability $O(n^{-B})$ for prescribed
$B$ and an explicit polynomial deviation bound. The proof follows the
log-ratio martingale of \cite[Lemma~6.7]{Lon26} with a sharper stopping
threshold and Freedman's inequality in place of the $L^2$ maximal bound; we
give the details because the rate is used essentially.

\begin{lemma}[Balanced depletion]\label{lem:balance}
Let $m$ targets be partitioned into $d$ categories of sizes $m_1,\dots,m_d$
with $m_b\ge c^*m^\vartheta$ for some $c^*>0$, $\vartheta\in(0,1]$. Let
targets be selected sequentially without replacement in an arbitrary
predictable adaptive order, and suppose that whenever $s$ targets remain,
every individual remaining target has conditional selection probability in
$[e^{-\beta}/s,e^{\beta}/s]$ with $0\le\beta\le\beta_0$. Fix
$\delta\in(0,1/4]$ and $B<\infty$. Then with probability at least
$1-O\bigl(d\,(m^{-B}+e^{-c(\log m_b)^2m_b^{1/2}/\delta})\bigr)$,
simultaneously for every category $b$ and every $r\le\delta m$,
\begin{equation}\label{eq:balance}
\Bigl|\,X_b(r)-\frac{rm_b}{m}\Bigr|
\;\le\;C_B\Bigl(\beta+\frac{\sqrt{\delta\log m}}{m_b^{1/4}}
+\frac{1}{m_b^{1/2}}\Bigr)m_b,
\end{equation}
where $X_b(r)$ is the number of category-$b$ targets among the first $r$
selections. The estimate remains valid when a fixed target is forbidden
throughout (cf.\ \cite[Lemma~6.7]{Lon26}).
\end{lemma}

\begin{proof}
Fix a category $b$ and write $N_r=m_b-X_b(r)$, $s_r=m-r$, following the proof
of \cite[Lemma~6.7]{Lon26}. Let $Y_{r+1}$ be the indicator that the
$(r+1)$-st selection lies in category $b$; summing the one-target bounds over
the $N_r$ remaining category-$b$ targets gives
$\Prob(Y_{r+1}=1\mid\cF_r)\in[e^{-\beta}N_r/s_r,e^{\beta}N_r/s_r]$
\cite[eq.~(6.6)]{Lon26}. Set $a=m_b^{1/2}$ and
$\tau=\inf\{r:N_r<a\text{ or }s_r<m/2\}$; since $r\le\delta m$ with
$\delta\le1/4$ gives $s_r\ge3m/4$, on our range $\tau$ is triggered only by
$\{N_r<a\}$. On $r<\tau$ put $Z_r=\log(N_r/s_r)$. As computed in
\cite[eq.~(6.7)]{Lon26},
\[
\bigl|\E[Z_{r+1}-Z_r\mid\cF_r]\bigr|
\;\le\;\frac{C\beta}{s_r}+C\Bigl(\frac{1}{s_r^2}+\frac{1}{s_rN_r}\Bigr),
\qquad
\Var(Z_{r+1}-Z_r\mid\cF_r)\;\le\;\frac{C}{s_rN_r},
\]
and $|Z_{r+1}-Z_r|\le C(s_r^{-1}+N_r^{-1})\le C'a^{-1}$ deterministically.
Let $M_r$ be the martingale of centered increments, stopped at $\tau$. For
$r\le\delta m\wedge\tau$: the total conditional variance is
\[
V\;\le\;\sum_{r\le\delta m}\frac{C}{s_rN_r}\;\le\;
\frac{C\delta m}{(m/2)\,a}\;=\;\frac{2C\delta}{a}\;=:\;v,
\]
the total drift is
$D\le C\beta\sum_{r\le\delta m}s_r^{-1}+CV\le C\beta\delta+C/a=O(\beta+1/a)$,
where we used
$\sum_{r\le\delta m}\frac1{m-r}=\log\frac{m}{m-\delta m}+O(m^{-1})=O(\delta)$,
and the maximal increment is $M\le C/a$. Freedman's inequality \cite{Fre75}
gives, for every $t>0$,
\[
\Prob\Bigl(\max_{r\le\delta m\wedge\tau}
\bigl|M_r\bigr|>t\Bigr)\;\le\;2\exp\Bigl(-\frac{t^2}{2(v+Mt/3)}\Bigr).
\]
Taking $t_1=C_B\bigl(\sqrt{\delta\log m}/a^{1/2}+(\log m)/a\bigr)$ with
$C_B$ large gives $t_1^2/(2(v+Mt_1/3))\ge B\log m$, hence the displayed
maximum exceeds $t_1$ with probability $\le2m^{-B}$. Taking
$t_2=\frac12\log(m_b/a)-D$ (which satisfies $t_2\ge\frac13\log m_b$ for
$\beta$ small and $m_b$ large), the event $\{\tau\le\delta m\}$ requires
$|M_\tau|\ge t_2$, since $Z_\tau\le\log(a/(m/2))$ while
$Z_0=\log(m_b/m)$ and the drift is at most $D$; hence
\[
\Prob(\tau\le\delta m)\;\le\;
2\exp\Bigl(-\frac{t_2^2}{2(v+Mt_2/3)}\Bigr)
\;\le\;2\exp\bigl(-c\,(\log m_b)^2\,a/\delta\bigr)
\]
because $v=2C\delta/a$ and $Mt_2\le C(\log m_b)/a=o(t_2^2)$.
On the complement of these two bad events, for every $r\le\delta m$,
\[
\log\frac{N_r/m_b}{s_r/m}\;=\;Z_r-\log\frac{m_b}{m}
\;\in\;\bigl[-D-t_1,\;D+t_1\bigr],
\]
hence, since $s_r/m\in[1/2,1]$ and $\beta\le\beta_0$ is small,
\[
\frac{N_r}{m_b}\;=\;\frac{s_r}{m}\,
\Bigl(1+O\bigl(\beta+a^{-1}+t_1\bigr)\Bigr),
\]
which is \eqref{eq:balance} after substituting
$X_b(r)=m_b-N_r$ and $a=m_b^{1/2}$ (the $(\log m)/a$ part of $t_1$ is
absorbed by the $m_b^{-1/2}$ term). Union over the $d$ categories gives the
stated probability. The forbidden-target variant is identical after replacing
$m$ by $m-1$, exactly as in \cite[Lemma~6.7]{Lon26}.
\end{proof}

\begin{remark}[Parameters; the balance event]\label{rem:balance-params}
We apply Lemma~\ref{lem:balance} to each row $a$ of the slab decomposition,
with targets = the $m_{n,a}$ rank positions of $Q_{n,a}$, categories = the
cells $Q_{n,a}\cap H_{n,b}$, selections = the orbit's targets in row $a$,
and $\beta=2\kappa h$; the one-target bounds $[e^{-\beta}/s,e^{\beta}/s]$
hold by (S3), including the forbidden-root variant under $\Prob^\circ$
\cite[Lemma~6.6]{Lon26}. On $\cG_n$ the cell sizes satisfy
$N_{n,ab}\ge cnh^2$ (concentration at means of order
$\E N_{n,ab}=\Theta(nh^2)$, which is part of $\cG_n$), so the hypothesis
$m_b\ge c^*m^\vartheta$ of Lemma~\ref{lem:balance} holds with
$m=m_{n,a}\ge cnh$ and $\vartheta=\log(nh^2)/\log(nh)>0$; the number of
cells per row is at most $d$ and the number of rows is $d$, so the union
bound multiplies the failure probability by $d^2=O(h^{-2}\log^2n)$, which is
$O(n^{-B})$ for the prescribed $B$ with our mesh choices. We write
\begin{equation}\label{eq:eps-bal}
\eps_n^{\mathrm{bal}}\;=\;
C\Bigl(\beta+(nh^2)^{-1/4}\sqrt{\log n}\Bigr),
\end{equation}
so that on the \emph{balance event} $\cB_n$,
$\Prob(\cB_n^c)=O(n^{-B})$, for every row $a$, every cell $(a,b)$ and every
prefix $r\le\delta n$ the residual counts satisfy
\begin{equation}\label{eq:resid-counts}
\widetilde N_{ab}(r)\;=\;
\frac{\widetilde m_a(r)}{m_{n,a}}\,N_{n,ab}\,
\bigl(1+O(\eps_n^{\mathrm{bal}})\bigr),
\qquad \widetilde m_a(r)=m_{n,a}-r_a,
\end{equation}
where $r_a$ is the number of exposed edges in row $a$.
\end{remark}

\subsection{The one-step slab kernel}
\begin{lemma}[One-step slab kernel]\label{lem:onestep}
On $\cB_n$, for every prefix $r\le\delta n$ with $\delta\le\delta_0(\kappa)$
small, every $a,b$, and every history of the $\Prob^\circ$
exploration,
\begin{equation}\label{eq:onestep}
\Prob^\circ\bigl(b_{r+1}=b\mid b_r=a,\ \cF_r,\ \cH_n\bigr)
\;=\;\bigl(1+O(\eta_n)\bigr)\,P_{n,ab},
\qquad \eta_n\;=\;C\bigl(\beta+\eps_n^{\mathrm{bal}}\bigr).
\end{equation}
\end{lemma}

\begin{proof}
By \eqref{eq:slab-transition} (the flat one-edge marginal (S3) summed over
the residual targets of cell $(a,b)$; the root-avoiding conditioning changes
only the forbidden-target normalization of (S3), whose relative effect is
$O(1/(nh))$, absorbed in $\eps_n^{\mathrm{bal}}$), the left side is
$e^{O(\beta)}\widetilde N_{ab}(r)/\widetilde m_a(r)$. By
\eqref{eq:resid-counts} this equals
$e^{O(\beta)}(1+O(\eps_n^{\mathrm{bal}}))N_{n,ab}/m_{n,a}
=(1+O(\beta+\eps_n^{\mathrm{bal}}))P_{n,ab}$.
\end{proof}

\subsection{Uniform mixing of the relative density}
Recall $u_k(b)=\Prob^\circ(b_k=b\mid\cH_n)/p_{n,b}$.

\begin{lemma}[Uniform mixing]\label{lem:mixing}
On $\cB_n$, with $q'=q_0+\eta_n$ (so $q'<1$ for $\eta_n$ small), for every
$k\ge1$ and every starting slab,
\begin{equation}\label{eq:mixing}
\bigl\|u_k-1\bigr\|_{L^2(p_n)}
\;\le\;(\kappa^A+1)\,(q')^{k-1}\;+\;\frac{\eta_n}{1-q'}
\;\le\;C(q')^k+C\eta_n.
\end{equation}
\end{lemma}

\begin{proof}
Write $\tilde P^{(r)}$ for the random (history-dependent) one-step kernel of
Lemma~\ref{lem:onestep}: $\tilde P^{(r)}_{ab}=(1+\eta_{ab}^{(r)})P_{n,ab}$
with $|\eta_{ab}^{(r)}|\le\eta_n$ on $\cB_n$. The density evolves by
$u_{r+1}(b)=\sum_a u_r(a)p_{n,a}\tilde P^{(r)}_{ab}/p_{n,b}$, i.e.
\[
u_{r+1}\;=\;P_n^*u_r+w_r,\qquad
w_r(b)\;=\;\sum_a u_r(a)\,p_{n,a}\,\eta_{ab}^{(r)}\,P_{n,ab}/p_{n,b}.
\]
Since all summands of $P_n^*u_r$ are nonnegative,
$|w_r|\le\eta_n P_n^*u_r$ pointwise, hence
$\|w_r\|_{L^2(p_n)}\le\eta_n\|P_n^*u_r\|_{L^2(p_n)}
\le\eta_n\|u_r\|_{L^2(p_n)}$ (Markov adjoints contract $L^2$). Moreover
$\E_{p_n}[w_r]=\E_{p_n}[u_{r+1}]-\E_{p_n}[P_n^*u_r]=1-1=0$; therefore
$P_n^*u_r-1=P_n^*(u_r-1)$ and, since $\E_{p_n}[u_r-1]=0$, the gap
\eqref{eq:gap} gives
\[
\|u_{r+1}-1\|_2\;\le\;q_0\|u_r-1\|_2+\eta_n(1+\|u_r-1\|_2)
\;\le\;q'\|u_r-1\|_2+\eta_n.
\]
After one step from any point start, the ratio bound \eqref{eq:ratio-bound}
gives $u_1(c)=(1+O(\eta_n))P_{n,b_0c}/p_{n,c}\in[\kappa^{-A}/2,2\kappa^A]$,
so $\|u_1-1\|_2\le\kappa^A+1$; the recursion
$a_{r+1}\le q'a_r+\eta_n$ from $r=1$ with $a_r=\|u_r-1\|_2$ yields
\eqref{eq:mixing}.
\end{proof}

\begin{lemma}[Visit-count fluctuations]\label{lem:visits}
On $\cB_n$, let $S_c(k)=\#\{0\le t\le k-1:b_t=c\}$ be the number of visits
of the orbit to slab $c$ among the first $k$ edges, $k\le\delta n$. Then
\begin{equation}\label{eq:visits}
\E^\circ S_c(k)\;\le\;k p_{n,c}\bigl(1+o(1)\bigr),
\qquad
\Var^\circ S_c(k)\;\le\;C k\,p_{n,c},
\end{equation}
and consequently, since $m_{n,c}=np_{n,c}$,
$\Prob^\circ\bigl(S_c(k)\ge m_{n,c}/2\bigr)
\le C\delta\,(nh)^{-1}$ for $\delta\le\delta_0$ small.
\end{lemma}

\begin{proof}
Write $S_c(k)=\sum_{t<k}\1\{b_t=c\}$. For the mean,
$\E^\circ\1\{b_t=c\}=p_{n,c}u_t(c)\le p_{n,c}(1+\|u_t-1\|_2p_{n,c}^{-1/2})$,
and summing \eqref{eq:mixing} over $1\le t<k$ gives
$\E^\circ S_c(k)\le kp_{n,c}+C\sqrt{p_{n,c}}
+Ck\eta_n\sqrt{p_{n,c}}=kp_{n,c}(1+o(1))$ since
$\sqrt{p_{n,c}}\le1$ and $k\eta_n\sqrt{p_{n,c}}\le
\delta n\eta_n=o(np_{n,c})$ with our mesh choices
($np_{n,c}\ge cnh$ and $\delta n\eta_n=o(nh)$ in both regimes below).
For the variance, let $j>t\ge1$: the difference $v_{j-t}$ of the densities
at time $j$ started from slabs $b_t$ and $b_0$ is the difference of two
$p_n$-densities, hence exactly mean-zero; each history-dependent kernel
$\tilde P^{(r)}=(1+\eta^{(r)})P_n$ contracts mean-zero vectors in
$L^2(p_n)$ by $q_0+\eta_n=q'$ (the gap \eqref{eq:gap} plus the entrywise
perturbation bound), so by \eqref{eq:ratio-bound} applied to the first step,
$\|v_{j-t}\|_2\le2\kappa^A(q')^{j-t-1}$, and
\[
\bigl|\Cov^\circ\bigl(\1\{b_t=c\},\1\{b_j=c\}\bigr)\bigr|
\;=\;\bigl|\E^\circ\bigl[\1\{b_t=c\}\,p_{n,c}\,v_{j-t}(c)\bigr]\bigr|
\;\le\;C\,p_{n,c}^{3/2}(q')^{j-t},
\]
using $\Prob^\circ(b_t=c)\le p_{n,c}(1+\|u_t-1\|_2p_{n,c}^{-1/2})
\le Cp_{n,c}$ for $t\ge1$ (the $t=0$ term contributes at most
$p_{n,c}^{1/2}\|u_0\|_2\le1$, affecting only the constant). Summing,
$\Var^\circ S_c(k)\le Ckp_{n,c}+Ckp_{n,c}^{3/2}/(1-q')\le Ckp_{n,c}$, and
Chebyshev with $\E^\circ S_c(k)\le m_{n,c}/3$ for $\delta\le\delta_0$
gives $\Prob^\circ(S_c(k)\ge m_{n,c}/2)\le
Ckp_{n,c}/m_{n,c}^2\le C\delta/(nh)$.
\end{proof}

\subsection{The root average in trace form}
The following computation is the finite-volume shadow of the trace
universality $r\lambda_r=\Tr(K^r)=1+O(q^r)$ of
\cite[Theorem~5.19]{Lon26}.

\begin{lemma}[Root-averaged return density]\label{lem:root-avg}
On $\cB_n$, for $1\le k\le\delta n$,
\begin{equation}\label{eq:root-avg}
R_k\;:=\;\frac1n\sum_{v=1}^n u_k^{(b(v))}\bigl(c(v)\bigr)
\;=\;\E^\circ\Bigl[\Tr\bigl(P_n\tilde P^{(1)}\cdots\tilde P^{(k)}\bigr)\Bigr]
\;=\;1+\Err_k,
\end{equation}
where $u_k^{(b)}$ is the density with start in slab $b$ and the expectation
is over the exploration history. The error satisfies:
\begin{enumerate}[label=(\roman*)]
\item for $k\le c_*\log n$ with mesh $h_k=(q')^k$ and
$c_*=\frac{1}{8|\log q'|}$: $\Err_k=O(q_0^k+k\eta_n)=o(1)$ as
$k\to\infty$, uniformly in $n$;
\item for $c_*\log n\le k\le\delta n$ with the fixed mesh $h=n^{-1/8}$:
$\Err_k=O(n^{-c_2})$.
\end{enumerate}
\end{lemma}

\begin{proof}
\emph{The trace identity.} Since
$\sum_v\1\{b(v)=b,c(v)=c\}=|H_{n,b}\cap Q_{n,c}|=N_{n,cb}$,
\[
R_k\;=\;\frac1n\sum_{b,c}N_{n,cb}\,u_k^{(b)}(c)
\;=\;\sum_{b,c}p_{n,c}P_{n,cb}\,\frac{\Prob^\circ(b_k=c\mid b_0=b)}{p_{n,c}}
\;=\;\E^\circ\sum_{b,c}P_{n,cb}\bigl(\tilde P^{(1)}\cdots\tilde
P^{(k)}\bigr)_{bc},
\]
and $\sum_bP_{n,cb}M_{bc}=(P_nM)_{cc}$, so the sum is
$\Tr(P_n\tilde P^{(1)}\cdots\tilde P^{(k)})$.

\emph{Regime (i): run expansion.} Write $P_n=1p_n+Q$ (so $Q1=0$ and
$p_nQ=0$ exactly, by Markovianity and stationarity) and
$\tilde P^{(s)}=P_n+E_s=1p_n+(Q+E_s)$ with $\|E_s\|_{L^2(p_n)\to L^2(p_n)}
\le\eta_n$ (entrywise bound as in Lemma~\ref{lem:mixing}). Expand the
product into stationary ($1p_n$) and transient ($Q+E_s$) factors. The
all-stationary term is $\Tr(1p_n)=1$ exactly. Every other term contains at
least one maximal run of transient factors; cyclically separating the trace
at the rank-one factors, each transient run $B$ contributes a scalar
$p_nB1$: since $p_nQ=0$, the run must start with an $E$, hence
$|p_nB1|\le\eta_n(q')^{r-1}$ for a run of length $r$ (each further factor
contributes $\|Q+E_s\|_{L^2_0\to L^2_0}\le q'$). Summing over at most
$k+1$ possible run positions and $t\ge1$ runs,
\[
\bigl|\Tr\bigl(P_n\tilde P^{(1)}\cdots\tilde P^{(k)}\bigr)-1\bigr|
\;\le\;\Tr(Q^{k+1})+\sum_{t\ge1}
\Bigl(\frac{(k+1)\eta_n}{1-q'}\Bigr)^t
\;=\;\Tr(Q^{k+1})+O(k\eta_n),
\]
for $k\eta_n$ small. Here the all-transient choice reproduces
$\Tr(Q^{k+1})$ plus terms with at least one $E$, which are absorbed in the
sum. By Weyl's inequality $\sum_i|\nu_i|^2\le\|Q\|_{\mathrm{HS}}^2
\le\|P_n\|_{\mathrm{HS}}^2\le\kappa^A$ (see \eqref{eq:hs}) applied to the
eigenvalues $\nu_i$ of $Q$ (all $|\nu_i|\le q_0$ by \eqref{eq:gap}),
$\Tr(Q^{k+1})=O(\kappa^Aq_0^{k-1})=O(q_0^k)$. In regime (i),
$\eta_n=C((q')^k+(n(q')^{2k})^{-1/4}\sqrt{\log n})
\le C((q')^k+n^{-3/16}\sqrt{\log n})$, so
$k\eta_n=o(1)$ for $k\le c_*\log n$ and $\Err_k=O(q_0^k+k\eta_n)$.

\emph{Regime (ii): half-split.} Write
$P_n\tilde P^{(1)}\cdots\tilde P^{(k)}
=M\bigl(\tilde P^{(\lfloor k/2\rfloor+1)}\cdots\tilde P^{(k)}\bigr)$ with
$M=P_n\tilde P^{(1)}\cdots\tilde P^{(\lfloor k/2\rfloor)}$. By
Lemma~\ref{lem:mixing}, every mean-zero vector contracts by
$\rho_k:=C(q')^{\lfloor k/2\rfloor}+C\eta_n/(1-q')$ under the second
factor, so $\tilde P^{(\lfloor k/2\rfloor+1)}\cdots\tilde P^{(k)}
=1p_n+\Delta$ with $\|\Delta\|_{L^2(p_n)\to L^2(p_n)}\le\rho_k$. Since the
second factor is Markov, $(1p_n+\Delta)1=1$, and
$\Tr(M(1p_n+\Delta))=\Tr((1p_n)M)+\Tr(M\Delta)=p_nM1+\Tr(M\Delta)
=1+\Tr(M\Delta)$ because $M1=1$ (Markov rows). By \eqref{eq:hs},
$\|M\|_{\mathrm{tr}}\le d^{1/2}\kappa^{A/2}$, hence
\[
|\Err_k|\;=\;|\Tr(M\Delta)|\;\le\;d^{1/2}\kappa^{A/2}\rho_k
\;\le\;C h^{-1/2}\sqrt{\log n}\,\bigl((q')^{k/2}+\eta_n\bigr)
\;=\;O(n^{-c_2})
\]
for $c_*\log n\le k\le\delta n$ with $h=n^{-1/8}$, since
$(q')^{c_*\log n/2}=n^{-1/16}$ and
$\eta_n=O(n^{-1/8}+n^{-3/16}\sqrt{\log n})$.
\end{proof}

\section{The closure hazard and the local cycle-length law}\label{sec:local}

\subsection{The hazard identity and its telescoping}
Recall from Proposition~\ref{prop:rootavoid} that
$S(k)=\#\{0\le t\le k-1:b_t=c(v)\}$ counts the closure opportunities and
$\lambda_j$ are the closure marginals \eqref{eq:hazard-bounds}.

\begin{lemma}[Hazard identity]\label{lem:hazard}
On $\cG_n$, for every root $v$ and every $k\ge1$,
\begin{equation}\label{eq:hazard}
\Prob\bigl(L_n(v)=k+1\mid\cH_n\bigr)
\;=\;\E^\circ\Bigl[\1\{b_k=c(v)\}
\prod_{j=0}^{S(k)-1}(1-\lambda_j)\,\lambda_{S(k)}\Bigr].
\end{equation}
\end{lemma}

\begin{proof}
By \eqref{eq:survival} applied at $k$ and at $k+1$,
\[
\Prob(L_n(v)=k+1\mid\cH_n)
\;=\;\E^\circ\Bigl[\prod_{j=0}^{S(k)-1}(1-\lambda_j)\Bigr]
-\E^\circ\Bigl[\prod_{j=0}^{S(k+1)-1}(1-\lambda_j)\Bigr].
\]
The edge exposed at step $k$ can close to $v$ only if the source $q_k$ lies
in the root's rank row, i.e.\ $b_k=c(v)$ (structural fact (S1): rank $v$ is
a target of row $c(v)$ only). On $\{b_k\ne c(v)\}$ one has
$S(k+1)=S(k)$ and the two products agree; on $\{b_k=c(v)\}$ one has
$S(k+1)=S(k)+1$ and the difference is the displayed integrand.
\end{proof}

\begin{lemma}[Telescoping]\label{lem:telescope}
On $\cG_n$, for every orbit history with $S(k)\le m_{n,c(v)}/2$,
\begin{equation}\label{eq:telescope}
\prod_{j=0}^{S(k)-1}(1-\lambda_j)\,\lambda_{S(k)}
\;=\;\frac{1+O(\beta)}{m_{n,c(v)}},
\end{equation}
and for every $S(k)\le m_{n,c(v)}-1$ the same quantity is at most
$C/m_{n,c(v)}$.
\end{lemma}

\begin{proof}
Write $m=m_{n,c(v)}$ and $\lambda_j=e^{\eps_j}/(m-j)$ with
$|\eps_j|\le\beta$ by \eqref{eq:hazard-bounds}. For the upper bound,
\[
1-\lambda_j\;\le\;1-\frac{e^{-\beta}}{m-j}
\;=\;\frac{m-j-1}{m-j}\Bigl(1+\frac{1-e^{-\beta}}{m-j-1}\Bigr)
\;\le\;\frac{m-j-1}{m-j}\,e^{\beta/(m-j-1)},
\]
so with $S=S(k)$,
\[
\prod_{j=0}^{S-1}(1-\lambda_j)\;\le\;
\frac{m-S}{m}\,e^{\beta\sum_{j<S}1/(m-j-1)}
\;\le\;\frac{m-S}{m}\,e^{\beta S/(m-S)},
\]
and multiplying by $\lambda_S\le e^{\beta}/(m-S)$ gives
$\le e^{\beta(1+S/(m-S))}/m$, which is $C/m$ for $S\le m-1$ and
$(1+O(\beta))/m$ for $S\le m/2$. The matching lower bound for $S\le m/2$ is
identical with $e^{+\beta}$ replaced by $e^{-\beta}$:
\[
1-\lambda_j\;\ge\;1-\frac{e^{\beta}}{m-j}
\;=\;\frac{m-j-1}{m-j}\Bigl(1-\frac{e^{\beta}-1}{m-j-1}\Bigr)
\;\ge\;\frac{m-j-1}{m-j}\,e^{-2\beta/(m-j-1)},
\]
giving $\prod(1-\lambda_j)\lambda_S\ge e^{-2\beta(1+S/(m-S))}/m$.
\end{proof}

\begin{lemma}[Closure probability given $\cH_n$]\label{lem:closure}
On $\cB_n$, for $\delta\le\delta_0(\kappa)$ small, every root $v$, and every
$k\le\delta n$,
\begin{equation}\label{eq:closure}
\Prob\bigl(L_n(v)=k+1\mid\cH_n\bigr)
\;=\;\frac{1+O(\beta)}{n}\,u_k\bigl(c(v)\bigr)
\;+\;O\bigl(\delta(nh)^{-2}\bigr).
\end{equation}
\end{lemma}

\begin{proof}
Split \eqref{eq:hazard} on the event $\{S(k)\le m_{n,c(v)}/2\}$. On that
event, Lemma~\ref{lem:telescope} makes the integrand equal
$(1+O(\beta))m_{n,c(v)}^{-1}\1\{b_k=c(v)\}$; off it, the same lemma bounds
the integrand by $Cm_{n,c(v)}^{-1}\1\{b_k=c(v)\}$. Hence
\[
\Prob(L_n(v)=k+1\mid\cH_n)
\;=\;\frac{1+O(\beta)}{m_{n,c(v)}}\,
\Prob^\circ\bigl(b_k=c(v)\mid\cH_n\bigr)
\;+\;O\Bigl(\frac{\Prob^\circ(S(k)\ge m_{n,c(v)}/2)}{m_{n,c(v)}}\Bigr).
\]
By Lemma~\ref{lem:visits},
\[
\frac{\Prob^\circ(S(k)\ge m_{n,c(v)}/2)}{m_{n,c(v)}}
\;\le\;\frac{C\delta}{nh\cdot cnh}
\;=\;O\bigl(\delta(nh)^{-2}\bigr),
\]
and $\Prob^\circ(b_k=c(v)\mid\cH_n)/m_{n,c(v)}
=p_{n,c(v)}u_k(c(v))/m_{n,c(v)}=u_k(c(v))/n$.
\end{proof}

\subsection{Proof of the local cycle-length law}
\begin{proof}[Proof of Theorem~\ref{thm:local}]
Average \eqref{eq:closure} over a uniform root $v=I_n$: by
Lemma~\ref{lem:root-avg},
\begin{equation}\label{eq:local-pf}
\Prob\bigl(L_n(I_n)=k+1\bigr)
\;=\;\frac{1+O(\beta)}{n}\,\bigl(1+\Err_k\bigr)+O\bigl(\delta(nh)^{-2}\bigr)
\;=\;\frac1n\Bigl(1+O\bigl(\beta+|\Err_k|\bigr)\Bigr)+o(1/n),
\end{equation}
since $(nh)^{-2}=o(1/n)$ in both mesh regimes. For $k\to\infty$ with
$k\le c_*\log n$ (mesh $h_k=(q')^k$), $\beta_k=2\kappa(q')^k$ and
$\Err_k=O(q_0^k+k\eta_n)=O(c_1^k)+O(n^{-c_2})$ with
$c_1=\max(q_0,q')^{1/2}$; for $c_*\log n\le k\le\delta n$ (mesh
$h=n^{-1/8}$), $\beta=2\kappa n^{-1/8}$ and $\Err_k=O(n^{-c_2})$. Hence for
every $r\in[r_n,\delta n]$ with $r_n\to\infty$,
\[
\Prob(L_n(I_n)=r)\;=\;\frac{1+o(1)}{n},
\qquad o(1)=O\bigl(c_1^{\,r}\bigr)+O(n^{-c_2}),
\]
uniformly in $r$, where the two mesh regimes are glued at
$c_*\log n$ (both give $O(n^{-c_2})$ there), and the $O(n^{-B})$ failure
probabilities of $\cG_n,\cB_n$ with $B\ge2$ are absorbed. This is
\eqref{eq:local-law} with the stated rate.
\end{proof}

\begin{proof}[Proof of Corollary~\ref{cor:EC}]
$\E C_r(\sigma_n)=\sum_{v=1}^n\Prob(L_n(v)=r)/r
=\frac{n}{r}\Prob(L_n(I_n)=r)=(1+o(1))/r$ by
Theorem~\ref{thm:local}.
\end{proof}

\begin{remark}[Why the two mesh regimes are needed]\label{rem:two-mesh}
The flatness error $O(\beta)=O(\kappa h)$ forces $h\to0$; the concentration
underlying $\cG_n$ and Lemma~\ref{lem:balance} forces $nh^2\to\infty$
polynomially; and the trace error of Lemma~\ref{lem:root-avg}(i) is
$O(q_0^k+k\eta_n)$, small only when the mesh is adapted to the scale $k$.
For $k\to\infty$ slowly (e.g.\ $k=\log\log n$) the scale-adaptive choice
$h_k=(q')^k$ makes all three errors vanish; beyond $k\asymp\log n$ the fixed
fine mesh $h=n^{-1/8}$ is more efficient and gives the polynomially small
error used in Section~\ref{sec:expectation}.
\end{remark}

\section{Iterated residual exploration and the joint local law}\label{sec:joint}

\subsection{Residual stability}
The framework of Section~\ref{sec:prelim} is stable under the deletion of
completely explored cycles, because a cycle intersects each row in as many
sources as targets.

\begin{lemma}[Residual stability]\label{lem:residual}
On $\cB_n$, let $\cE$ be the exploration of the cycles of uniform roots
$V_1,\dots,V_{j-1}$, of total length $K=o(n)$, all distinct. Let $R$ be the
residual vertex set, $|R|=n-K$, and let $\sigma_n^R$ be the residual
permutation (the unexposed matching). Then:
\begin{enumerate}[label=(\roman*)]
\item \textbf{Square rows.} $|R\cap H_{n,a}|=|R\cap Q_{n,a}|=:m_a^R$ for
every $a$, and $m_a^R=m_{n,a}-r_a$ with
$r_a=(K/n)m_{n,a}(1+O(\eps_n^{\mathrm{bal}}))$; similarly
$N_{ab}^R=N_{n,ab}(1-K/n)(1+O(\eps_n^{\mathrm{bal}}))$.
\item \textbf{Kernel perturbation.} The residual compressed kernel
$P_{ab}^R=N_{ab}^R/m_a^R$ satisfies
$P_{ab}^R=(1+O(\eps_n^{\mathrm{bal}}))P_{n,ab}$, its stationary law $p^R$
satisfies $p_a^R=(1+o(1))p_{n,a}$, and on $\cB_n$,
\begin{equation}\label{eq:resid-gap}
\|P^Rv\|_{L^2(p^R)}\;\le\;(q_0+o(1))\,\|v\|_{L^2(p^R)}
\qquad\text{whenever }\E_{p^R}[v]=0.
\end{equation}
\item \textbf{Flatness and one-edge marginals} persist with the same
$\beta=2\kappa h$ \cite[Lemmas~6.5, 6.6]{Lon26}, as does the root-avoiding
change of measure for a residual root \cite[Lemma~6.9]{Lon26}.
\end{enumerate}
Consequently every conclusion of Sections~\ref{sec:orbit} and
\ref{sec:local} holds for $\sigma_n^R$ with $n$ replaced by $n-K$ and
unchanged exponents.
\end{lemma}

\begin{proof}
(i) Each explored cycle $C$ satisfies $|C\cap H_{n,a}|=|C\cap Q_{n,a}|$ by
(S1) (the complete-cycle deletion identity of \cite[Section~6.2]{Lon26}),
so rows stay square; the balanced form of $r_a$ and $N_{ab}^R$ is
\eqref{eq:resid-counts} applied to the total prefix $K\le\delta n$.
(ii) The entrywise bound is \eqref{eq:resid-counts} divided through by
$m_a^R/m_{n,a}$; it gives $\|P^R-P_n\|_{L^2(p_n)\to L^2(p_n)}
=O(\eps_n^{\mathrm{bal}})$ since
$\|(P^R-P_n)v\|_2^2=\sum_a p_{n,a}(\sum_b(P^R-P_n)_{ab}v_b)^2
\le O(\eps^2)\|P_n|v|\|_2^2$. Also
$\|p^R/p_n-1\|_\infty=O(K/n+\eps_n^{\mathrm{bal}})=o(1)$ by (i), so
mean-zero functions in $L^2(p^R)$ are $o(\|v\|_2)$-close to mean-zero in
$L^2(p_n)$, and \eqref{eq:resid-gap} follows from \eqref{eq:gap} with
$q_0+O(\eps_n^{\mathrm{bal}})+o(1)=q_0+o(1)$.
(iii) is hereditary flatness (S3).
\end{proof}

\subsection{The joint local law}
\begin{proof}[Proof of Theorem~\ref{thm:joint-local}]
Fix $\ell$ and $k_1,\dots,k_\ell\to\infty$ with $K_\ell=\sum_j k_j=o(n)$.
We prove by induction on $j=1,\dots,\ell$ that
\begin{equation}\label{eq:induct}
\Prob\Bigl(\bigcap_{i\le j}\{L_n(V_i)=k_i,\ V_i\notin
C(V_1)\cup\cdots\cup C(V_{i-1})\}\Bigr)
\;=\;\frac{1+o(1)}{n^j},
\end{equation}
with $o(1)$ uniform in the lengths; the event in \eqref{eq:induct} at
$j=\ell$ is exactly the event of Theorem~\ref{thm:joint-local} (given
$V_i\notin\bigcup_{i'<i}C(V_{i'})$ for all $i$, the cycles are distinct).
The case $j=1$ is Theorem~\ref{thm:local}:
$\Prob(L_n(V_1)=k_1)=(1+o(1))/n$. For the step, condition on the
exploration $\cE$ of the first $j-1$ cycles (total length
$K_{j-1}=\sum_{i<j}k_i=o(n)$), on the balance event. Then
\[
\Prob\bigl(V_j\notin\textstyle\bigcup_{i<j}C(V_i)\mid\cE\bigr)
\;=\;\frac{n-K_{j-1}}{n-j+1}\;=\;1-\frac{K_{j-1}}{n}+O(n^{-1}),
\]
and, given $V_j$ in the residual set, $V_j$ is uniform there, so
Lemma~\ref{lem:residual} applied to $\sigma_n^R$ with the local law
\eqref{eq:local-pf} gives
\[
\Prob\bigl(L_n(V_j)=k_j\mid\cE,\ V_j\in R\bigr)
\;=\;\frac{1+o(1)}{n-K_{j-1}},
\]
with $o(1)$ uniform: the mesh and all exponents of
Sections~\ref{sec:orbit}--\ref{sec:local} depend only on $\kappa$ and the
mesh, and $n-K_{j-1}=n(1+o(1))$. Multiplying,
\[
\Prob\bigl(V_j\notin\textstyle\bigcup_{i<j}C(V_i),\ L_n(V_j)=k_j
\mid\cE\bigr)
\;=\;\frac{1+o(1)}{n-K_{j-1}}\cdot
\Bigl(1-\frac{K_{j-1}}{n}\Bigr)\;=\;\frac{1+o(1)}{n},
\]
the cancellation $(n-K_{j-1})^{-1}(1-K_{j-1}/n)=1/n$ being exact up to
$O(n^{-1})$. Multiplying into \eqref{eq:induct} at $j-1$ gives the claim at
$j$; the accumulated error is $(1+o(1))^\ell=1+o(1)$ since $\ell$ is fixed,
and all failure events are $O(\ell n^{-B})=O(n^{-B})$.
\end{proof}

\begin{remark}[The size-biased bookkeeping]\label{rem:size-bias}
Theorem~\ref{thm:joint-local} controls counts of \emph{cycle tuples}. For
target lengths $k_1,\dots,k_\ell$ (not necessarily distinct), let
$\mathfrak{C}_n(k_1,\dots,k_\ell)$ be the number of ordered $\ell$-tuples of
distinct cycles of $\sigma_n$ with the $j$-th cycle of length $k_j$. Marking
each cycle by a uniform vertex gives
$\E\mathfrak{C}_n(k_1,\dots,k_\ell)
=\frac{n^\ell(1+o(1))}{\prod_j k_j}\,
\Prob(\text{event of }\eqref{eq:joint-local})
=\frac{1+o(1)}{\prod_j k_j}$,
since each ordered cycle tuple corresponds to exactly $\prod_j k_j$ ordered
vertex tuples on distinct cycles. This is the mesoscopic analogue of the
exact uniform-permutation identity
$\E\mathfrak{C}_n(k_1,\dots,k_\ell)=\prod_j k_j^{-1}$
\cite{ABT92}.
\end{remark}

\section{The mesoscopic Poisson process: proof of Theorem~\ref{thm:ppp}}\label{sec:ppp}

Let $a_n\to\infty$, $a_n=o(n)$, and let $J_1,\dots,J_m$ be pairwise disjoint
intervals $J_j=[u_j,v_j]\Subset(0,\infty)$. Write
$\cN_n(J)=\#\{\text{cycles }C:|C|/a_n\in J\}$.

\begin{lemma}[Factorial moments of mesoscopic counts]\label{lem:meso-fm}
For every fixed $\ell_1,\dots,\ell_m\ge0$,
\begin{equation}\label{eq:meso-fm}
\E\prod_{j=1}^m\bigl(\cN_n(J_j)\bigr)_{\ell_j}
\;\longrightarrow\;\prod_{j=1}^m\Bigl(\log\frac{v_j}{u_j}\Bigr)^{\ell_j}
\qquad(n\to\infty).
\end{equation}
\end{lemma}

\begin{proof}
The falling factorial $(\cN_n(J))_\ell$ counts ordered $\ell$-tuples of
distinct cycles with rescaled lengths in $J$. Hence, with
$\ell=\sum_j\ell_j$ and Remark~\ref{rem:size-bias},
\[
\E\prod_j\bigl(\cN_n(J_j)\bigr)_{\ell_j}
\;=\;\sum_{(k_{j,i})}\E\,\mathfrak{C}_n\bigl((k_{j,i})\bigr)
\;=\;\bigl(1+o(1)\bigr)\sum_{(k_{j,i})}\prod_{j,i}\frac1{k_{j,i}},
\]
where the sum runs over all $\ell$-tuples of lengths with
$k_{j,i}\in a_nJ_j=[u_ja_n,v_ja_n]$, and the joint law applies because each
$k_{j,i}\ge u_ja_n\to\infty$ and
$\sum_{j,i}k_{j,i}\le\ell\,(\max_jv_j)a_n=o(n)$, so the uniformity of
Theorem~\ref{thm:joint-local} covers all summands. The error $o(1)$ is
uniform over the $(a_n)^\ell$-many length tuples (each application carries
the error $O(c_1^{u_ja_n}+n^{-c_2})$ of Theorem~\ref{thm:local} per
coordinate, uniform over the range), so
\[
\E\prod_j\bigl(\cN_n(J_j)\bigr)_{\ell_j}
\;=\;\bigl(1+o(1)\bigr)\prod_{j=1}^m
\Bigl(\sum_{u_ja_n\le k\le v_ja_n}\frac1k\Bigr)^{\ell_j}
\;=\;\bigl(1+o(1)\bigr)\prod_{j=1}^m
\Bigl(\log\frac{v_j}{u_j}+o(1)\Bigr)^{\ell_j},
\]
which is \eqref{eq:meso-fm}.
\end{proof}

\begin{proof}[Proof of Theorem~\ref{thm:ppp}]
The right side of \eqref{eq:meso-fm} is the joint factorial moment of
independent $\Poisson(\log(v_j/u_j))$ variables. Poisson distributions are
determined by their moments (Carleman's condition
$\sum_k\mu_k^{-1/2k}=\infty$ holds for Poisson moments), and convergence of
all joint factorial moments implies convergence of all joint ordinary
moments (the two triangular transforms are equivalent), hence joint
convergence in distribution to the independent Poisson vector by the method
of moments (see e.g.\ \cite{BHJ92}, where the same argument establishes
Poisson limits for cycle counts). This is exactly
\eqref{eq:ppp-vector}. The equivalent formulation as convergence of
$\sum_i\delta_{L_{n,i}/a_n}$ to the Poisson point process with intensity
$dx/x$ on sets bounded away from $0$ and $\infty$ follows since the
factorial moment measures converge:
$\E\sum^*\prod_j\delta_{L_j/a_n}(J_j)\to\prod_j\int_{J_j}dx/x$ for all
disjoint boxes, which characterizes the Poisson process with that intensity
\cite{ABT92,Pit06}. This resolves both parts of
\cite[Conjecture~13.1]{Lon26}.
\end{proof}

\section{The expected number of cycles: proof of Theorem~\ref{thm:expectation}}\label{sec:expectation}

Fix $\delta\in(0,\delta_0]$ with $\delta_0=\delta_0(\kappa)$ from
Section~\ref{sec:local} and a large fixed $K$. Decompose
\begin{equation}\label{eq:four-ranges}
\E N_n^{\cyc}\;=\;\sum_{r=1}^n \E C_r(\sigma_n)
\;=\;\Bigl(\sum_{r\le K}+\sum_{K<r\le c_*\log n}
+\sum_{c_*\log n<r\le\delta n}+\sum_{r>\delta n}\Bigr)\E C_r(\sigma_n)
\;=:\;S_1+S_2+S_3+S_4.
\end{equation}

\emph{Range 1 (fixed $r$).} By Proposition~\ref{prop:short}(i),
$S_1\to\sum_{r=1}^K\lambda_r$.

\emph{Range 2 ($K<r\le c_*\log n$).} By Theorem~\ref{thm:local} with its
rate, uniformly in $r$,
\[
\E C_r(\sigma_n)\;=\;\frac{n}{r}\,\Prob(L_n(I_n)=r)
\;=\;\frac1r\Bigl(1+O\bigl(c_1^{\,r}\bigr)+O(n^{-c_2})\Bigr),
\]
hence
\[
S_2\;=\;\sum_{K<r\le c_*\log n}\frac1r
\;+\;O\Bigl(\sum_{r>K}\frac{c_1^{\,r}}{r}\Bigr)
\;+\;O\Bigl(\frac{\log\log n}{n^{c_2}}\Bigr)
\;=\;\sum_{K<r\le c_*\log n}\frac1r+O\bigl(c_1^{\,K}\bigr)+o(1).
\]

\emph{Range 3 ($c_*\log n<r\le\delta n$).} Theorem~\ref{thm:local} gives
$\E C_r=(1+O(n^{-c_2}))/r$ uniformly, hence
\[
S_3\;=\;\sum_{c_*\log n<r\le\delta n}\frac1r
\;+\;O\Bigl(\frac{\log n}{n^{c_2}}\Bigr)
\;=\;\sum_{c_*\log n<r\le\delta n}\frac1r+o(1).
\]

\emph{Range 4 ($r>\delta n$).} By Proposition~\ref{prop:macro},
$S_4=\E\#\{i:L_{n,i}\ge\delta n\}\to\log(1/\delta)$ as $n\to\infty$ with
$\delta$ fixed, while
$\sum_{r>\delta n}1/r=\log(1/\delta)+o(1)$.

Assembling,
\[
\E N_n^{\cyc}
\;=\;\sum_{r\le K}\lambda_r+\sum_{K<r\le\delta n}\frac1r
+\log\frac1\delta+O\bigl(c_1^{\,K}\bigr)+o(1)
\;=\;\log n+\Bigl(\sum_{r\le K}\lambda_r-\log K\Bigr)
+O\bigl(c_1^{\,K}\bigr)+o(1),
\]
where the $\log\delta$ terms cancel exactly:
$\sum_{K<r\le\delta n}1/r=\log(\delta n/K)+o(1)$. Writing
$\sum_{r\le K}\lambda_r-\log K
=\sum_{r\le K}(\lambda_r-\frac1r)+(H_K-\log K)
=\sum_{r\le K}(\lambda_r-\frac1r)+\gE+o_K(1)$,
we obtain, for every fixed $K$,
\[
\limsup_{n\to\infty}\Bigl|\E N_n^{\cyc}-\log n-\kappa_f\Bigr|
\;\le\;\Bigl|\sum_{r>K}\Bigl(\lambda_r-\frac1r\Bigr)\Bigr|
+O\bigl(c_1^{\,K}\bigr),
\]
and the right side tends to $0$ as $K\to\infty$ by the absolute convergence
\eqref{eq:trace-gap}. Hence $\E N_n^{\cyc}=\log n+\kappa_f+o(1)$.
\qed

\begin{remark}\label{rem:uniform-in-delta}
The proof shows the decomposition is uniform in the choice of the split
point $\delta\le\delta_0$: the mesoscopic local law contributes the full
harmonic sum over $(K,\delta n]$ with summable error, and the macroscopic
count supplies exactly the missing $\log(1/\delta)$. The constant
$\kappa_f$ is therefore precisely the accumulation of the short-cycle trace
corrections $\lambda_r-1/r=O(q^r/r)$ over all scales.
\end{remark}

\section{The central limit theorem: proof of Theorem~\ref{thm:clt}}\label{sec:clt}

Fix $\delta\in(0,\delta_0]$ and $K$, and put
\[
N_n^{\mathrm{meso}}\;=\;\sum_{K<r\le\delta n}C_r(\sigma_n),
\qquad
\mu_n\;=\;\sum_{K<r\le\delta n}\frac1r
\;=\;\log\frac{\delta n}{K}+o(1).
\]

\begin{lemma}[Mesoscopic factorial moments]\label{lem:clt-fm}
For every fixed $\ell\ge1$,
\begin{equation}\label{eq:clt-fm}
\E\bigl(N_n^{\mathrm{meso}}\bigr)_\ell
\;=\;\mu_n^\ell\bigl(1+o(1)\bigr).
\end{equation}
\end{lemma}

\begin{proof}
As in Lemma~\ref{lem:meso-fm},
\[
\E\bigl(N_n^{\mathrm{meso}}\bigr)_\ell
\;=\;\sum_{K<r_1,\dots,r_\ell\le\delta n}
\E\,\mathfrak{C}_n(r_1,\dots,r_\ell)
\;=\;\bigl(1+o(1)\bigr)
\sum_{K<r_1,\dots,r_\ell\le\delta n}\prod_{j=1}^\ell\frac1{r_j}
\;=\;\bigl(1+o(1)\bigr)\mu_n^\ell,
\]
where Theorem~\ref{thm:joint-local} applies uniformly over the sum since
each $r_j>K\to\infty$ (the error is uniform over the range by the rate in
Theorem~\ref{thm:local}; equal lengths among the $r_j$ are permitted since
the cycles in a tuple are distinct but their lengths need not be) and
$\sum_j r_j\le\ell\delta n=o(n)$.
\end{proof}

\begin{lemma}[Mesoscopic asymptotic normality]\label{lem:meso-normal}
\[
\frac{N_n^{\mathrm{meso}}-\mu_n}{\sqrt{\mu_n}}\;\Rightarrow\;N(0,1).
\]
\end{lemma}

\begin{proof}
The factorial moments \eqref{eq:clt-fm} agree to relative error $o(1)$ with
those of $\Poisson(\mu_n)$. Ordinary moments are fixed triangular
combinations of factorial ones, $x^j=\sum_{i\le j}S(j,i)(x)_i$ with $S(j,i)$
the Stirling numbers of the second kind, so
$\E(N_n^{\mathrm{meso}})^j=\E[\Poisson(\mu_n)^j](1+o(1))$ for every fixed
$j$. Hence every fixed moment of
$(N_n^{\mathrm{meso}}-\mu_n)/\sqrt{\mu_n}$ agrees asymptotically with the
corresponding moment of the standardized Poisson, which converges to the
Gaussian moment ($0$ in odd order, $(j-1)!!$ in even order) as
$\mu_n\to\infty$; see e.g.\ \cite{BHJ92} for the standard computation.
Since the Gaussian law is moment-determinate,
$(N_n^{\mathrm{meso}}-\mu_n)/\sqrt{\mu_n}\Rightarrow N(0,1)$ by the method
of moments.
\end{proof}

\begin{proof}[Proof of Theorem~\ref{thm:clt}]
Write
\[
N_n^{\cyc}\;=\;N_n^{\mathrm{meso}}+N_n^{\mathrm{short}}+N_n^{\mathrm{macro}},
\qquad
N_n^{\mathrm{short}}=\sum_{r\le K}C_r(\sigma_n),\quad
N_n^{\mathrm{macro}}=\sum_{r>\delta n}C_r(\sigma_n).
\]
For fixed $K$, $N_n^{\mathrm{short}}\Rightarrow\sum_{r\le K}Z_r$ by
Proposition~\ref{prop:short}(i), hence is tight; for fixed $\delta$,
$N_n^{\mathrm{macro}}\Rightarrow\Poisson(\log(1/\delta))$ by
Proposition~\ref{prop:macro}, hence is tight. Therefore
$R_n=N_n^{\mathrm{short}}+N_n^{\mathrm{macro}}=O_p(1)$ and, since
$\mu_n=\log n+O(1)\to\infty$,
\[
\frac{N_n^{\cyc}-\mu_n}{\sqrt{\mu_n}}
\;=\;\frac{N_n^{\mathrm{meso}}-\mu_n}{\sqrt{\mu_n}}+\frac{R_n}{\sqrt{\mu_n}}
\;\Rightarrow\;N(0,1)
\]
by Lemma~\ref{lem:meso-normal} and Slutsky's theorem. Finally
$\E N_n^{\cyc}=\log n+\kappa_f+o(1)=\mu_n+O(1)$ by
Theorem~\ref{thm:expectation}, and
$\sqrt{\mu_n}=\sqrt{\log n}(1+o(1))$, so centering at $\E N_n^{\cyc}$ and
normalizing by $\sqrt{\log n}$ does not change the limit:
\[
\frac{N_n^{\cyc}-\E N_n^{\cyc}}{\sqrt{\log n}}\;\Rightarrow\;N(0,1).
\]
This resolves \cite[Conjecture~13.2]{Lon26}.
\end{proof}

\begin{remark}[Variance]\label{rem:variance}
The same computation gives
$\Var N_n^{\mathrm{meso}}=\mu_n(1+o(1))$, so
$\Var N_n^{\cyc}=\log n+o(\log n)$: the variance is asymptotic to the mean,
exactly as for the uniform permutation, where
$\Var N^{\cyc}=\log n+\gE-\pi^2/6+o(1)$ \cite{ABT92}. The trace
corrections $\lambda_r-1/r$ affect the $O(1)$ terms of the mean and
variance but not the leading $\log n$ of either.
\end{remark}
\section{The remaining conjectures: reductions and
obstructions}\label{sec:remaining}

\subsection{Conjecture A.1: reduction to a mixed-scale
estimate}\label{sec:A1}

\begin{estimate}[Mixed-scale local estimate]\label{est:mixed}
Let $(k_{n,j})_{j\ge1}$ be an admissible set of scales for the slab
decomposition of \cite[Section~6]{Lon26} in the sense of
\cite[Definition~8.1]{Lon26}. For each fixed $\ell\ge1$, distinct fixed
scales $s_1,\dots,s_\ell\in(0,\infty)$, and distinct mesoscopic lengths
$k_{n,1},\dots,k_{n,\ell'}$ with $k_{n,j}\to\infty$ and
$k_{n,j}=o(n)$,
\begin{equation}\label{eq:mixed}
\lim_{n\to\infty}\Big|
\Prob\big(\text{$L_{n,i}=k_{n,i}$ for $i\le\ell'$ and
$L_{n,\ell'+j}/n\to s_j$ for $j\le\ell$}\big)
-\frac{1+o(1)}{\prod_{i=1}^{\ell'}k_{n,i}}\cdot
\Prob\big(L_{n,\ell'+j}/n\to s_j,\ j\le\ell\big)\Big|=0,
\end{equation}
where the second factor is the macroscopic marginal of
\cite[Theorem~6.1]{Lon26} (the $\PD(1)$ finite-dimensional law), and
the convergence of the macroscopic coordinates is in the sense of the
finite-dimensional distributions established there.
\end{estimate}

\begin{theorem}[Reduction]\label{thm:A1-reduction}
Estimate~\ref{est:mixed} implies Conjecture~A.1 of \cite[Appendix~A]{Lon26}:
under \eqref{eq:mixed}, the point process of all rescaled inter-point
distances of the slab decomposition converges to the Poisson--Dirichlet
fixed point \cite[Theorem~10.1]{Lon26} jointly with independence of the
bulk and endpoint contributions.
\end{theorem}

\begin{proof}
By \cite[Lemma~8.2]{Lon26} the total variation distance between the
inter-point distance process and the point process
$\sum_{v}\delta_{L_n(v)/n}$ is $O(nh)$. Convergence of the full process
to the fixed point, and the asserted independence, are therefore
consequences of the joint convergence of the rescaled cycle lengths.
The joint finite-dimensional convergence in turn follows from
\eqref{eq:mixed} by summation over the mixed event and tightness of the
macroscopic vector (the coordinatewise bound
$\E\#\{i:L_{n,i}\ge\delta n\}\to\log(1/\delta)$ of \cite[Proposition~6.14]{Lon26},
uniformly over admissible scales); the factorization
\eqref{eq:mixed} identifies the limit as the product of the
scale-invariant mesoscopic law and the macroscopic $\PD(1)$ law, which
is precisely the independence asserted in Conjecture~A.1.
\end{proof}

\begin{remark}[Status of Estimate~\ref{est:mixed}]\label{rem:est-status}
Estimate~\ref{est:mixed} is stated for the \emph{full} Poisson--Dirichlet
model of \cite[Appendix~A]{Lon26}, in which macroscopic and mesoscopic
scales interact. We do \emph{not} prove it here. The two extreme
marginals are known: the purely mesoscopic marginal is
Theorem~\ref{thm:joint-local}, and the purely macroscopic marginal is
\cite[Theorem~10.1]{Lon26} (via the marginalization of
\cite[Theorem~6.1]{Lon26}). The missing ingredient is a single-run
exclusion argument in the regime where a macroscopic orbit (length
$\ge\delta n$) depletes a slab by a positive proportion before the
mesoscopic orbits close; our balanced-depletion lemma requires the
depletion proportion $\delta$ to be small relative to the flatness
scale, and the perturbation expansion of Lemma~\ref{lem:residual}
requires $k/n\to0$. A proof along the present lines would require
either a truncated-martingale closure argument valid for $k/n$
bounded away from $0$, or a direct coupling of the macroscopic orbit
to its limiting scale-invariant dynamics (the GEM allocation of
\cite[Section~9]{Lon26}) before conditioning the mesoscopic orbits.
We record Estimate~\ref{est:mixed} as the precise missing statement.
\end{remark}

\subsection{Conjecture 13.3: the comparable-weights
obstruction}\label{sec:133}

\begin{proposition}[Reduction]\label{prop:133}
Conjecture~13.3 of \cite[Section~13]{Lon26} (endpoint-singular
$\PD(1)$) follows from the local law Theorem~\ref{thm:local} together
with the local law at the singular scale $k=\Theta(n)$.
\end{proposition}

\begin{proof}
Conjecture~13.3 asserts that the endpoint-singular weight profiles
$\theta_{n,i}=\Theta(i^{-\alpha})$ ($\alpha>0$) produce the same
$\PD(1)$ macroscopic law. The proof of \cite[Theorem~6.1]{Lon26} from
the local law proceeds through: (i) the mesoscopic mass bound (any
fixed $\eps$-mass in lengths $\le\delta n$ contributes $\le2\eps$
mass), which follows from Theorem~\ref{thm:local} and
Corollary~\ref{cor:EC} as in Section~\ref{sec:expectation}; (ii)
tightness, from the macroscopic moment bound; (iii) identification of
the limit via the induced partition, which requires the local law at
scales $k=\Theta(n)$ (the endpoint scale), where the comparable-weights
hypothesis \cite[eq.~(2.1)]{Lon26} fails.
\end{proof}

\begin{remark}[The obstruction]\label{rem:obstruction}
The comparable-weights hypothesis is used in \cite{Lon26} and in the
present paper in three places: the spectral gap
(Lemma~\ref{prop:good}, via \cite[Lemmas~6.2--6.3]{Lon26}), the
hereditary flatness bound $\beta=2\kappa h$
(\cite[Lemmas~6.5--6.6]{Lon26}), and the ratio bound
\cite[eq.~(6.4)]{Lon26}. For the endpoint-singular profiles of
Conjecture~13.3, $\kappa=\kappa_n\to\infty$ (indeed
$\theta_{n,1}/\theta_{n,n}=\Theta(n^{\alpha})$), so all three inputs
degenerate: the gap bound $q_0$ is no longer $<1$ uniformly, the
flatness window $e^{\pm2\kappa_nh}$ is void unless
$h=o(1/\kappa_n)$, which forces slabs of width $o(n/\kappa_n)$ and
breaks the count concentrations, and the ratio bound degenerates.
A proof of Conjecture~13.3 along the present lines would require a
weighted (non-uniform) slab decomposition adapted to the profile, with
a gap estimate for a non-uniformly-ergodic compressed chain; this is a
genuine extension of the framework, not a consequence of it. We
therefore leave Conjecture~13.3 open, with
Proposition~\ref{prop:133} recording the exact reduction.
\end{remark}

\subsection{Conjecture A.2}\label{sec:A2}

Conjecture~A.2 (the law of the largest inter-point distance) requires
the upper tail of the cycle-length distribution at scale
$\Theta(n/\log n)$, a regime where our local law error $O(c_1^{\,r})$
is not $o(1/r)$. The Hammersley-type coupling of \cite[Section~11]{Lon26}
identifies the limiting largest-gap law conditionally on a joint
decoupling estimate at that scale, which is beyond the reach of the
present methods (for the same reason as in
Remark~\ref{rem:obstruction}: the scale $\Theta(n/\log n)$ is
macroscopic in the sense $k/n\not\to0$). We record the conjecture as
open.

\appendix

\section{Numerical verification}\label{app:numerics}

Theorems~\ref{thm:expectation} and \ref{thm:clt} were checked against
direct simulation of the Luce permutation in the exact sampled-profile
model \eqref{eq:profile} with $f(x)=1+x$ on $[0,1]$, i.e.\
$\theta_{n,i}=1+i/n$. With $A(t)=\int_0^1(1-e^{-f(u)t})\,du$ and
$g=f\circ A$, the substitution $x=A(s)$, $y=A(t)$ in
\eqref{eq:permuton} gives the time-domain kernel
$J(s,t)=g(s)e^{-g(s)t}$ on $(0,\infty)^2$, which is trace-equivalent to
$K$: $\Tr(K^r)=\Tr(J^r)$. Discretizing $J$ on $(0,30]$ and computing
direct matrix-power traces gives
\[
\lambda_1=\Tr(J)=0.8520,\qquad
\Tr(J^2)=1.0266,\qquad
\Tr(J^3)=0.9957,\qquad
\Tr(J^4)=1.0006,
\]
(the value of $\Tr(J)$ agrees with the position-domain quadrature
$\int_0^1\rho(x,x)\,dx=0.8520$ to four digits), whence, with
$\lambda_r=\Tr(J^r)/r$,
\[
\kappa_f\;=\;\gE+\sum_{r\ge1}\Bigl(\lambda_r-\frac1r\Bigr)
\;=\;0.4410,
\]
the tail $\sum_{r\ge5}$ contributing less than $10^{-3}$.
Theorem~\ref{thm:expectation} therefore predicts
$\E N_n^{\cyc}=\log n+0.4410+o(1)$. Direct simulation of the Luce
permutation ($6000$ trials per $n$, cycles counted by pointer doubling)
gives:
\[
\begin{array}{c|cccc}
n & 500 & 1000 & 2000 & 4000\\\hline
\E N_n^{\cyc} & 6.64 & 7.37 & 8.04 & 8.74\\
\log n+\kappa_f & 6.66 & 7.35 & 8.04 & 8.74\\
\Var N_n^{\cyc}/\log n & 0.81 & 0.83 & 0.85 & 0.85
\end{array}
\]
The short-cycle counts agree with the Poisson means $\lambda_r$ of
\cite[Theorem~5.1]{Lon26}: at $n=4000$ the empirical $\E C_r$ for
$r=1,2,3,4,5,6,8,10$ are $0.855,0.510,0.329,0.252,0.195,0.164,0.129,
0.101$, against $\lambda_r=0.852,0.513,0.332,0.250,0.200,0.167,0.125,
0.100$. The mesoscopic scale-invariance of Corollary~\ref{cor:EC} is
visible in dyadic windows: at $n=4000$ the mean number of cycles with
length in $[2^j,2^{j+1})$ is $0.69$--$0.73$ for $j=3,\dots,10$, against
the limiting $\log 2=0.693$. The variance ratio
$\Var N_n^{\cyc}/\log n$ increases toward the limit $1$ of
Remark~\ref{rem:variance}.

\begin{thebibliography}{99}

\bibitem[Lon26]{Lon26}
C.~D. Long,
\emph{Flux and fragmentation in non-uniform random permutations},
manuscript under referee revision, 2026.

\bibitem[ABT92]{ABT92}
R.~Arratia, A.~D. Barbour, and S.~Tavar\'e,
\emph{Poisson process approximations for the Ewens sampling formula},
Ann. Appl. Probab. \textbf{2} (1992), 519--535.

\bibitem[BHJ92]{BHJ92}
A.~D. Barbour, L.~Holst, and S.~Janson,
\emph{Poisson Approximation},
Oxford University Press, 1992.

\bibitem[BCD25]{BCD25}
J.~Borga, S.~Chatterjee, and P.~Diaconis,
\emph{Random permutations with comparable weights: stationarity and mixing},
preprint, 2025.

\bibitem[CDK24]{CDK24}
S.~Chatterjee, P.~Diaconis, and D.~Kane,
\emph{The Luce choice model and random permutations with prescribed
weights}, preprint, 2024.

\bibitem[DZ99]{DZ99}
A.~Dembo and O.~Zeitouni,
\emph{Large Deviations Techniques and Applications}, 2nd ed.,
Springer, 1999.

\bibitem[Fre75]{Fre75}
D.~A. Freedman,
\emph{On tail probabilities for martingales},
Ann. Probab. \textbf{3} (1975), 100--118.

\bibitem[Luc59]{Luc59}
R.~D. Luce,
\emph{Individual Choice Behavior: A Theoretical Analysis},
Wiley, 1959.

\bibitem[Petrov]{Petrov}
V.~V. Petrov,
\emph{Limit Theorems of Probability Theory: Sequences of Independent
Random Variables}, Oxford University Press, 1995.

\bibitem[Pit06]{Pit06}
J.~Pitman,
\emph{Combinatorial Stochastic Processes},
Lecture Notes in Mathematics \textbf{1875}, Springer, 2006.

\bibitem[Sj\"o23]{Sjo23}
J.~Sj\"ostrand,
\emph{Permutations with comparable weights: spectral asymptotics and
determinants}, preprint, 2023.

\end{thebibliography}

\end{document}
